\chapter{2010年上海卷（秋理）}

\section{填空题}

\begin{problem}
不等式 \(\frac{2 - x}{x + 4} > 0\) 的解集为\fillinblank{}.
\end{problem}

\begin{answer}
\((-4,2)\)
\end{answer}

\begin{solution}
分式有意义的条件是 \(x\neq-4\). 令分子、分母的零点分别为 \(x=2\) 和 \(x=-4\)，在数轴上将其分成 \((-\infty,-4)\)、\((-4,2)\)、\((2,+\infty)\) 三段. 取 \(x=-5\) 时分式为负，取 \(x=0\) 时分式为正，取 \(x=3\) 时分式为负，因此不等式的解集为 \((-4,2)\).
\end{solution}

\begin{problem}
若复数 \(z = 1 - 2\i\) （\(\i\) 为虚数单位），则 \(z \cdot \overline{z} + z =\fillinblank{}\).
\end{problem}

\begin{answer}
\(6-2\i\)
\end{answer}

\begin{solution}
因为 \(z=1-2\i\)，所以 \(\overline{z}=1+2\i\). 于是 \(z\overline{z}=(1-2\i)(1+2\i)=1+4=5\)，从而 \(z\overline{z}+z=5+(1-2\i)=6-2\i\).
\end{solution}

\begin{problem}
动点 \(P\) 到点 \(F(2,0)\) 的距离与它到直线 \(x + 2 = 0\) 的距离相等，则点 \(P\) 的轨迹方程为\fillinblank{}.
\end{problem}

\begin{answer}
\(y^2=8x\)
\end{answer}

\begin{solution}
设动点 \(P(x,y)\). 由抛物线的定义，点 \(P\) 到焦点 \(F(2,0)\) 的距离等于点 \(P\) 到准线 \(x+2=0\) 的距离，因此
\[
\sqrt{(x-2)^2+y^2}=|x+2|
\]
两边平方得 \((x-2)^2+y^2=(x+2)^2\)，整理为 \(y^2=8x\). 这就是点 \(P\) 的轨迹方程.
\end{solution}

\begin{problem}
行列式 \(\begin{vmatrix} \cos \frac{\pi}{3} & \sin \frac{\pi}{6} \\ \sin \frac{\pi}{3} & \cos \frac{\pi}{6} \end{vmatrix}\) 的值是\fillinblank{}.
\end{problem}

\begin{answer}
0
\end{answer}

\begin{solution}
代入特殊角的三角函数值，行列式为
\[
\begin{vmatrix}
\frac12&\frac12\\
\frac{\sqrt{3}}2&\frac{\sqrt{3}}2
\end{vmatrix}
=\frac12\cdot\frac{\sqrt{3}}2-\frac12\cdot\frac{\sqrt{3}}2=0
\]
所以填 \(0\).
\end{solution}

\begin{problem}
圆 \(C: x^{2} + y^{2} - 2x - 4y + 4 = 0\) 的圆心到直线 \(l: 3x + 4y + 4 = 0\) 的距离 \(d =\fillinblank{}\).
\end{problem}

\begin{answer}
3
\end{answer}

\begin{solution}
将圆的方程配方，得
\[
(x-1)^2+(y-2)^2=1
\]
所以圆心为 \(C_0(1,2)\). 由点到直线的距离公式，圆心到直线 \(3x+4y+4=0\) 的距离为
\[
d=\frac{|3\cdot1+4\cdot2+4|}{\sqrt{3^2+4^2}}=\frac{15}{5}=3
\]
\end{solution}

\begin{problem}
随机变量 \(\xi\) 的概率分布由下表给出：

\begin{center}
\begin{tabular}{c|cccc}
\(x\) & 7 & 8 & 9 & 10 \\
\hline
\(P(\xi = x)\) & 0.3 & 0.35 & 0.2 & 0.15
\end{tabular}
\end{center}

则该随机变量 \(\xi\) 的均值是\fillinblank{}.
\end{problem}

\begin{answer}
8.2
\end{answer}

\begin{solution}
表中概率之和为 \(0.3+0.35+0.2+0.15=1\)，故均值按离散型随机变量的定义计算为
\[
E(\xi)=7\times0.3+8\times0.35+9\times0.2+10\times0.15
=2.1+2.8+1.8+1.5=8.2
\]
\end{solution}

\begin{problem}
2010 年上海世博会园区每天 9:00 开园，20:00 停止入园. 在下边的框图中，\(S\) 表示上海世博会官方网站在每个整点报道的入园总人数，\(a\) 表示整点报道前 1 个小时内入园人数，则空白的执行框内应填入\fillinblank{}.

\texfigure[max width=0.40\linewidth]{img_repaint/2010/shanghai_science/q7_fig1.tex}
\end{problem}

\begin{answer}
\(S\leftarrow S+a\)
\end{answer}

\begin{solution}
在循环开始时，\(S\) 已保存此前各整点报道的入园人数总和. 读入整点报道前 \(1\) 小时新增的入园人数 \(a\) 后，应把这部分人数累加到总人数中，因此执行框应填写 \(S\leftarrow S+a\).
\end{solution}

\begin{problem}
对任意不等于 1 的正数 \(a\)，函数 \(f(x)=\log_{a}(x+3)\) 的反函数的图象都经过点 \(P\)，则点 \(P\) 的坐标为\fillinblank{}.
\end{problem}

\begin{answer}
\((0,-2)\)
\end{answer}

\begin{solution}
设反函数图象上的横坐标为 \(x\)，纵坐标为 \(y\). 由 \(y=\log_a(x+3)\) 交换 \(x\)，\(y\) 得 \(x=\log_a(y+3)\)，所以反函数为 \(y=a^x-3\). 代入 \(x=0\)，得 \(y=a^0-3=-2\)，此值与正数 \(a\neq1\) 无关. 于是所有这些反函数图象都经过 \(P(0,-2)\).
\end{solution}

\begin{problem}
从一副混合后的扑克牌 （\(52\text{ 张}\)） 中，随机抽取 \(1\) 张，事件 \(A\) 为“抽得红桃 \(K\)”，事件 \(B\) 为“抽得为黑桃”，则概率 \(P(A \cup B) =\fillinblank{}\). （结果用最简分数表示）
\end{problem}

\begin{answer}
\(\frac{7}{26}\)
\end{answer}

\begin{solution}
一副扑克牌共有 \(52\) 张. 事件 \(A\) 只有红桃 \(K\) 这 \(1\) 张牌，事件 \(B\) 包含 \(13\) 张黑桃. 两事件不相交，因此
\[
P(A\cup B)=P(A)+P(B)=\frac{1}{52}+\frac{13}{52}=\frac{14}{52}=\frac{7}{26}
\]
\end{solution}

\begin{problem}
在 \(n\) 行 \(n\) 列矩阵 \(\begin{pmatrix}1&2&3&\cdots&n-2&n-1&n\\2&3&4&\cdots&n-1&n&1\\3&4&5&\cdots&n&1&2\\\vdots&\ddots&\vdots&\ddots&\vdots&\cdots&\vdots\\n&1&2&\cdots&n-3&n-2&n-1\end{pmatrix}\) 中，记位于第 \(i\) 行第 \(j\) 列的数为 \(a_{ij}\)（\(i\)，\(j=1\)，\(2\)，\(\cdots\)，\(n\)）. 当 \(n = 9\) 时，\(a_{11} + a_{22} + a_{33} + \cdots + a_{99} =\fillinblank{}\).
\end{problem}

\begin{answer}
45
\end{answer}

\begin{solution}
矩阵每一行都把上一行的数整体加 \(1\)，超过 \(n\) 后从 \(1\) 重新开始. 当 \(n=9\) 时，对角线上的元素依次为
\(a_{11}=1\)，\(a_{22}=3\)，\(a_{33}=5\)，\(a_{44}=7\)，\(a_{55}=9\)
\(a_{66}=2\)，\(a_{77}=4\)，\(a_{88}=6\)，\(a_{99}=8\).
所以
\[
a_{11}+a_{22}+\cdots+a_{99}=1+3+5+7+9+2+4+6+8=45
\]
\end{solution}

\begin{problem}
将直线 \(l_1: nx + y - n = 0\)，\(l_2: x + ny - n = 0\) （\(n \in \N^*\)，\(n \geqslant 2\)），\(x\) 轴，\(y\) 轴围成的封闭区域的面积记为 \(S_n\)，则 \(\displaystyle\lim_{n \to \infty} S_n =\fillinblank{}\).
\end{problem}

\begin{answer}
1
\end{answer}

\begin{solution}
两直线与坐标轴的交点分别为：\(l_1\) 与 \(x\) 轴交于 \((1,0)\)，与 \(y\) 轴交于 \((0,n)\)；\(l_2\) 与 \(x\) 轴交于 \((n,0)\)，与 \(y\) 轴交于 \((0,1)\). 取被两条坐标轴、两直线围成的靠近原点的封闭区域.

联立两直线方程
\(nx+y=n\)，\(x+ny=n\)
相减得 \(x=y\)，代回得交点 \(Q\left(\frac{n}{n+1},\frac{n}{n+1}\right)\). 该区域可分成三角形 \(O(0,0)\)，\((0,1)\)，\(Q\) 和三角形 \(O\)，\(Q\)，\((1,0)\)，故
\[
S_n=\frac12\cdot1\cdot\frac{n}{n+1}+\frac12\cdot1\cdot\frac{n}{n+1}=\frac{n}{n+1}
\]
因此
\[
\lim_{n\to\infty}S_n=\lim_{n\to\infty}\frac{n}{n+1}=1
\]
\end{solution}

\begin{problem}
如图所示，在边长为 4 的正方形纸片 \(ABCD\) 中，\(AC\) 与 \(BD\) 相交于点 \(O\)，剪去 \(\triangle AOB\). 将剩余部分沿 \(OC\)，\(OD\) 折叠，使 \(OA\)，\(OB\) 重合，则以 \(A(B)\)，\(C\)，\(D\)，\(O\) 为顶点的四面体的体积是\fillinblank{}.

\bitmapfigure[max width=0.30\linewidth]{img/2010/shanghai_science/q12_fig1.png}
\end{problem}

\begin{answer}
\(\frac{8\sqrt{2}}{3}\)
\end{answer}

\begin{solution}
折叠后，三角形 \(COD\) 作为四面体的底面，原来的点 \(A\)、\(B\) 重合为顶点 \(X=A(B)\). 取空间直角坐标系，使
\(O=(0,0,0)\)、\(C=(2,2,0)\)、\(D=(-2,2,0)\). 由折叠保持长度，得
\(XO=AO=2\sqrt2\)，\(XC=BC=4\)，\(XD=AD=4\). 由 \(XC=XD\)，可设 \(X=(0,y,z)\). 于是
\(y^2+z^2=8\)，\(4+(y-2)^2+z^2=16\).
两式相减得 \(y=0\)，再由第一式得 \(z^2=8\). 因而四面体的高为 \(2\sqrt2\).

底面三角形 \(COD\) 的底 \(CD=4\)，且 \(O\) 到 \(CD\) 的距离为 \(2\)，所以底面积为 \(4\). 因此四面体体积为
\[
V=\frac13\cdot4\cdot2\sqrt2=\frac{8\sqrt2}{3}
\]
\end{solution}

\begin{problem}
如图所示，直线 \(x = 2\) 与双曲线 \(\Gamma: \frac{x^2}{4} - y^2 = 1\) 的渐近线交于 \(E_1\)，\(E_2\) 两点，记 \(\overrightarrow{OE_1} = \bs{e}_1\)，\(\overrightarrow{OE_2} = \bs{e}_2\)，任取双曲线 \(\Gamma\) 上的点 \(P\)，若 \(\overrightarrow{OP} = a\bs{e}_1 + b\bs{e}_2\)（\(a\)，\(b \in \R\)），则 \(a\)，\(b\) 满足的一个等式是\fillinblank{}.

\bitmapfigure[max width=0.55\linewidth]{img/2010/shanghai_science/q13_fig1.png}
\end{problem}

\begin{answer}
\(4ab=1\)
\end{answer}

\begin{solution}
双曲线的渐近线为 \(y=\frac{x}{2}\) 和 \(y=-\frac{x}{2}\). 与直线 \(x=2\) 的交点可取 \(E_1=(2,1)\)、\(E_2=(2,-1)\)，所以
\(\bs e_1=(2,1)\)，\(\bs e_2=(2,-1)\).
由 \(\overrightarrow{OP}=a\bs e_1+b\bs e_2\)，得 \(P=(2a+2b,a-b)\). 代入双曲线方程，得到
\[
\frac{(2a+2b)^2}{4}-(a-b)^2=(a+b)^2-(a-b)^2=4ab=1
\]
因此 \(a\)，\(b\) 满足的等式为 \(4ab=1\).
\end{solution}

\begin{problem}
从集合 \(U = \{a, b, c, d\}\) 的子集中选出 4 个不同的子集，需同时满足以下两个条件：

\begin{circlelist}
\item \(\varnothing\)，\(U\) 都要选出；
\item 对选出的任意两个子集 \(A\) 和 \(B\)，必有 \(A \subseteq B\) 或 \(A \supseteq B\).
\end{circlelist}

那么，共有\fillinblank{} 种不同的选法.
\end{problem}

\begin{answer}
36
\end{answer}

\begin{solution}
因为必须选出 \(\varnothing\) 和 \(U\)，其余两个子集必须是两个非空真子集. 条件（2）要求所有选出的子集两两可比较，因此其余两个子集 \(A\)，\(B\) 必须满足 \(A\subset B\).

按 \(A\)、\(B\) 的元素个数分类计数：
当 \(|A|=1\)，\(|B|=2\) 时，有 \(\mathrm C_4^1\mathrm C_3^1=12\) 种；当 \(|A|=1\)，\(|B|=3\) 时，有 \(\mathrm C_4^1\mathrm C_3^2=12\) 种；当 \(|A|=2\)，\(|B|=3\) 时，有 \(\mathrm C_4^2\mathrm C_2^1=12\) 种.
三类互不相交，故不同选法共有 \(12+12+12=36\) 种.
\end{solution}

\section{单选题}

\begin{problem}
\(x = 2k\pi + \frac{\pi}{4}\) （\(k \in \Z\)） 是“\(\tan x = 1\)”成立的（\quad）

\choices
    {充分不必要条件}
    {必要不充分条件}
    {充要条件}
    {既不充分也不必要条件}

\begin{answer}
A
\end{answer}

\begin{solution}
设条件 \(p\)：\(x=2k\pi+\frac{\pi}{4}\)，\(k\in\Z\)，条件 \(q\)：\(\tan x=1\). 若 \(p\) 成立，则
\(\tan x=\tan\left(2k\pi+\frac{\pi}{4}\right)=1\)，所以 \(p\) 是 \(q\) 的充分条件.

但 \(q\) 的一般解为 \(x=k\pi+\frac{\pi}{4}\)，例如 \(x=\frac{5\pi}{4}\) 满足 \(q\)，却不能写成 \(2k\pi+\frac{\pi}{4}\). 所以 \(p\) 不是 \(q\) 的必要条件，故选 A，即充分不必要条件.
\end{solution}
\end{problem}

\begin{problem}
直线 \(l\) 的参数方程是 \(\begin{cases}
x = 1 + 2t,\\
y = 2 - t,
\end{cases}\) （\(t \in \R\)），则 \(l\) 的方向向量 \(\bs{d}\) 可以是（\quad）

\choices
    {\((1,2)\)}
    {\((2,1)\)}
    {\((-2,1)\)}
    {\((1,-2)\)}

\begin{answer}
C
\end{answer}

\begin{solution}
直线参数方程中 \(t\) 的系数给出一个方向向量，因此直线 \(l\) 的方向向量可取 \(\bs d=(2,-1)\). 方向向量的任意非零倍数都可以表示同一方向，而选项 C 中
\((-2,1)=-(2,-1)\)，所以选 C.
\end{solution}
\end{problem}

\begin{problem}
若 \(x_0\) 是方程 \(\left(\frac{1}{2}\right)^x = x^{\frac{1}{3}}\) 的解，则 \(x_0\) 属于区间（\quad）

\choices
    {\(\left(\frac{2}{3},1\right)\)}
    {\(\left(\frac{1}{2}, \frac{2}{3}\right)\)}
    {\(\left(\frac{1}{3}, \frac{1}{2}\right)\)}
    {\(\left(0, \frac{1}{3}\right)\)}

\begin{answer}
C
\end{answer}

\begin{solution}
左边 \(\left(\frac12\right)^x\) 始终为正，所以方程的解必须满足 \(x>0\). 对 \(x>0\) 两边取自然对数，得
\(x\ln\frac12=\frac13\ln x\)，\(\text{即 }x\cdot 8^x=1\).
函数 \(g(x)=x8^x\) 在 \(x>0\) 上严格递增，因为
\(g'(x)=8^x(1+x\ln8)>0\)，所以方程至多有一个正根.

又
\(g\left(\frac13\right)=\frac13\cdot2=\frac23<1\)，\(g\left(\frac12\right)=\frac12\cdot\sqrt8=\sqrt2>1\).
由连续性，方程的根属于 \(\left(\frac13,\frac12\right)\)，故选 C.
\end{solution}
\end{problem}

\begin{problem}
某人要作一个三角形，要求它的三条高的长度分别是 \(\frac{1}{13}\)，\(\frac{1}{11}\)，\(\frac{1}{5}\)，则此人将（\quad）

\choices
    {不能作出满足要求的三角形}
    {作出一个锐角三角形}
    {作出一个直角三角形}
    {作出一个钝角三角形}

\begin{answer}
D
\end{answer}

\begin{solution}
设三角形面积为 \(\Delta\)，三条高分别为 \(h_a\)，\(h_b\)，\(h_c\)，对应边分别为 \(a\)，\(b\)，\(c\). 由
\(\Delta=\frac12ah_a=\frac12bh_b=\frac12ch_c\)，得边长之比为
\[
a:b:c=\frac1{h_a}:\frac1{h_b}:\frac1{h_c}=13:11:5
\]
因为 \(11+5>13\)，这三个数能作成三角形. 其中最长边为 \(13\)，且
\[
13^2=169>11^2+5^2=146
\]
由余弦定理的逆定理，最长边所对的角为钝角，所以应作出一个钝角三角形，选 D.
\end{solution}
\end{problem}

\section{解答题}

\begin{problem}
已知 \(0 < x < \frac{\pi}{2}\)，化简： \(\lg \left(\cos x \cdot \tan x + 1 - 2 \sin^2 \frac{x}{2}\right) + \lg \left[ \sqrt{2} \cos \left( x - \frac{\pi}{4} \right) \right] - \lg (1 + \sin 2x)\).
\end{problem}

\begin{solution}
由 \(0 < x < \frac{\pi}{2}\)，得 \(\sin x>0\)、\(\cos x>0\)，从而 \(\sin x+\cos x>0\). 又
\(\cos x\cdot\tan x+1-2\sin^2\frac{x}{2}=\sin x+1-(1-\cos x)=\sin x+\cos x\)，
\(\sqrt{2}\cos\left(x-\frac{\pi}{4}\right)=\sin x+\cos x\)，且
\(1+\sin 2x=1+2\sin x\cos x=(\sin x+\cos x)^2\).
因此原式为 \(\lg(\sin x+\cos x)+\lg(\sin x+\cos x)-\lg\left((\sin x+\cos x)^2\right)=0\).
\end{solution}

\begin{problem}
已知数列 \(\{a_{n}\}\) 的前 \(n\) 项和为 \(S_{n}\)，且 \(S_{n} = n - 5a_{n} - 85\)，\(n \in \N^*\).

\begin{enumerate}
\item 证明：\(\{a_{n}-1\}\) 是等比数列；
\item 求数列 \(\{S_n\}\) 的通项公式，并指出 \(n\) 为何值时，\(S_n\) 取得最小值，并说明理由.
\end{enumerate}
\end{problem}

\begin{solution}
\begin{enumerate}
\item 当 \(n=1\) 时，\(S_1=a_1\)，代入已知关系得 \(a_1=1-5a_1-85\)，所以 \(a_1=-14\)，从而 \(a_1-1=-15\).

当 \(n\geqslant2\) 时，\(a_n=S_n-S_{n-1}\)，故
\(a_n=(n-5a_n-85)-\bigl(n-1-5a_{n-1}-85\bigr)=1-5a_n+5a_{n-1}\).
整理得 \(a_n-1=\frac{5}{6}(a_{n-1}-1)\). 结合首项 \(a_1-1=-15\neq0\)，可知 \(\{a_n-1\}\) 是首项为 \(-15\)、公比为 \(\frac{5}{6}\) 的等比数列.

\item 由上一问，\(a_n-1=-15\left(\frac{5}{6}\right)^{n-1}\)，即
\(a_n=1-15\left(\frac{5}{6}\right)^{n-1}\). 代入已知关系，得
\(S_n=n-5a_n-85=n+75\left(\frac{5}{6}\right)^{n-1}-90\).

相邻两项之差为 \(S_{n+1}-S_n=1-\frac{25}{2}\left(\frac{5}{6}\right)^{n-1}\). 直接比较幂的大小：
\[
\left(\frac{5}{6}\right)^{14}<\frac{2}{25}<\left(\frac{5}{6}\right)^{13}
\]
其中左侧等价于 \(25\cdot5^{14}<2\cdot6^{14}\)，即 \(152587890625<156728328192\)，右侧等价于 \(25\cdot5^{13}>2\cdot6^{13}\)，即 \(30517578125>26121388032\).
由于 \(\frac{5}{6}<1\)，当 \(n\leqslant14\) 时，\(\left(\frac{5}{6}\right)^{n-1}\geqslant\left(\frac{5}{6}\right)^{13}>\frac{2}{25}\)，所以 \(S_{n+1}-S_n<0\)；当 \(n\geqslant15\) 时，\(\left(\frac{5}{6}\right)^{n-1}\leqslant\left(\frac{5}{6}\right)^{14}<\frac{2}{25}\)，所以 \(S_{n+1}-S_n>0\). 因而 \(\{S_n\}\) 先减后增，且 \(n=15\) 时 \(S_n\) 取得最小值.
\end{enumerate}
\end{solution}

\begin{problem}
如图所示，为了制作一个圆柱形灯笼，先要制作 \(4\) 个全等的矩形骨架，总计耗用 \(9.6\text{ 米}\) 铁丝. 骨架将圆柱底面 \(8\) 等分，再用 \(S\text{ 平方米}\) 塑料片制成圆柱的侧面和下底面 （不安装上底面）.

\begin{enumerate}
\item 当圆柱底面半径 \(r\) 取何值时，\(S\) 取得最大值？并求出该最大值 （结果精确到 \(0.01\text{ 平方米}\)）；
    \item 在灯笼内，以矩形骨架的顶点为端点，安装一些霓虹灯. 当灯笼底面半径为 \(0.3\text{ 米}\) 时，求图中两根直线型霓虹灯 \(A_{1}B_{3}\)，\(A_{3}B_{5}\) 所在异面直线所成角的大小 （结果用反三角函数表示）.
\end{enumerate}

\bitmapfigure[max width=0.40\linewidth]{img/2010/shanghai_science/q21_fig1.png}
\end{problem}

\begin{solution}
\begin{enumerate}
\item 设圆柱的母线长为 \(h\). 每个矩形骨架的两条水平边都是圆柱底面的直径，长度均为 \(2r\)，两条竖直边长度均为 \(h\). 因而 \(4\) 个矩形骨架所耗铁丝总长满足 \(4\cdot2(2r+h)=9.6\)，即 \(h=1.2-2r\)，且 \(0<r<0.6\).

塑料片覆盖圆柱侧面和一个底面，所以 \(S=2\pi rh+\pi r^2=2\pi r(1.2-2r)+\pi r^2=-3\pi(r-0.4)^2+0.48\pi\). 因此，当 \(r=0.4\text{ 米}\) 时，\(S\) 取得最大值，最大值为 \(0.48\pi\text{ 平方米}\)，精确到 \(0.01\text{ 平方米}\) 为 \(1.51\text{ 平方米}\).

\item 当 \(r=0.3\text{ 米}\) 时，\(h=1.2-2\times0.3=0.6\). 建立空间直角坐标系，使圆柱底面在 \(z=0\) 平面上，圆心为原点，并取 \(A_1=(-r,0,0)\)、\(A_3=(0,-r,0)\)、\(A_5=(r,0,0)\). 于是 \(B_3=(0,-r,h)\)、\(B_5=(r,0,h)\)，从而 \(\overrightarrow{A_1B_3}=(r,-r,h)\)、\(\overrightarrow{A_3B_5}=(r,r,h)\).

设两根异面直线所成的锐角为 \(\theta\)，则
\[
\cos\theta=\frac{\overrightarrow{A_1B_3}\cdot\overrightarrow{A_3B_5}}{\left|\overrightarrow{A_1B_3}\right|\left|\overrightarrow{A_3B_5}\right|}=\frac{h^2}{2r^2+h^2}=\frac{0.6^2}{2\times0.3^2+0.6^2}=\frac{2}{3}
\]
所以两根异面直线所成的角为 \(\arccos\frac{2}{3}\).
\end{enumerate}
\end{solution}

\begin{problem}
若实数 \(x\)，\(y\)，\(m\) 满足 \(|x - m| > |y - m|\)，则称 \(x\) 比 \(y\) 远离 \(m\).

\begin{enumerate}
\item 若 \(x^{2} - 1\) 比 1 远离 0，求 \(x\) 的取值范围；
\item 对任意两个不相等的正数 \(a\)，\(b\)，证明： \(a^3 + b^3\) 比 \(a^2b + ab^2\) 远离 \(2ab\sqrt{ab}\)；
\item 已知函数 \(f(x)\) 的定义域 \(D = \left\{ x \mid x \neq \frac{k\pi}{2} + \frac{\pi}{4}, k \in \Z, x \in \R \right\}\). 任取 \(x \in D\)，\(f(x)\) 等于 \(\sin x\) 和 \(\cos x\) 中远离 \(0\) 的那个值. 写出函数 \(f(x)\) 的解析式，并指出它的基本性质 （结论不要求证明）.
\end{enumerate}
\end{problem}

\begin{solution}
\begin{enumerate}
\item 由“远离”的定义，需满足 \(\left|(x^2-1)-0\right|>|1-0|\)，即 \(\left|x^2-1\right|>1\). 因为 \(x^2\geqslant0\)，不可能有 \(x^2-1<-1\)，所以只能有 \(x^2-1>1\)，即 \(x^2>2\). 因此 \(x\in(-\infty,-\sqrt{2})\cup(\sqrt{2},+\infty)\).

\item 令 \(q=\sqrt{ab}>0\)，\(t=\frac{a+b}{q}\). 由 \(a\neq b\) 及均值不等式，\(a+b>2\sqrt{ab}\)，故 \(t>2\). 记 \(A=a^3+b^3\)、\(B=a^2b+ab^2\)、\(C=2ab\sqrt{ab}\). 由 \(a+b=qt\)、\(ab=q^2\)，得 \(A=q^3(t^3-3t)\)、\(B=q^3t\)、\(C=2q^3\). 因而 \(A-C=q^3(t^3-3t-2)=q^3(t-2)(t+1)^2>0\)，而 \(B-C=q^3(t-2)>0\). 又因为 \((t+1)^2>1\)，所以 \(A-C>B-C>0\)，从而
\(\left|a^3+b^3-2ab\sqrt{ab}\right|>\left|a^2b+ab^2-2ab\sqrt{ab}\right|\). 这正说明 \(a^3+b^3\) 比 \(a^2b+ab^2\) 远离 \(2ab\sqrt{ab}\).

\item 对于 \(x\in D\)，有 \(\left|\sin x\right|>\left|\cos x\right|\iff\sin^2x>\cos^2x\iff\cos2x<0\). 因此，当 \(x\in\left(k\pi+\frac{\pi}{4},k\pi+\frac{3\pi}{4}\right)\) 时，较远的值为 \(\sin x\)；当 \(x\in\left(k\pi-\frac{\pi}{4},k\pi+\frac{\pi}{4}\right)\) 时，较远的值为 \(\cos x\). 所以
\[
f(x)=
\begin{cases}
\sin x,&x\in\left(k\pi+\frac{\pi}{4},k\pi+\frac{3\pi}{4}\right),\\
\cos x,&x\in\left(k\pi-\frac{\pi}{4},k\pi+\frac{\pi}{4}\right),
\end{cases}
\qquad k\in\Z
\]
由正弦、余弦的周期性，\(f(x+2\pi)=f(x)\). 函数的最大值为 \(1\)，且恰在 \(x=2k\pi\) 或 \(x=2k\pi+\frac{\pi}{2}\) 时取得；这些最大值点组成的集合的最小正平移量为 \(2\pi\)，故 \(f\) 的最小正周期为 \(2\pi\). 此外，\(f(0)=1\neq0\)，所以 \(f\) 不是奇函数；又 \(f\left(-\frac{\pi}{2}\right)=-1\neq1=f\left(\frac{\pi}{2}\right)\)，所以 \(f\) 不是偶函数.

由于 \(x\in D\) 时 \(|\sin x|\neq|\cos x|\)，较远的那个值的绝对值严格大于 \(\frac{\sqrt{2}}{2}\)，且 \(\sin x\)、\(\cos x\) 分别在相应区间内取得从 \(\frac{\sqrt{2}}{2}\) 到 \(1\) 以及从 \(-1\) 到 \(-\frac{\sqrt{2}}{2}\) 的所有值，其中端点 \(\pm\frac{\sqrt{2}}{2}\) 因不属于定义域而不能取到. 当 \(x=2k\pi\) 或 \(x=2k\pi+\frac{\pi}{2}\) 时，\(f(x)=1\)；当 \(x=(2k+1)\pi\) 或 \(x=2k\pi+\frac{3\pi}{2}\) 时，\(f(x)=-1\)，其中 \(k\in\Z\). 因而值域为 \(\left[-1,-\frac{\sqrt{2}}{2}\right)\cup\left(\frac{\sqrt{2}}{2},1\right]\).

由上述分段式及正、余弦函数的单调性，\(f(x)\) 在
\[
\begin{gathered}
\left(2k\pi-\frac{\pi}{4},2k\pi\right),\quad \left(2k\pi+\frac{\pi}{4},2k\pi+\frac{\pi}{2}\right),\\
\left((2k+1)\pi,(2k+1)\pi+\frac{\pi}{4}\right),\quad \left((2k+1)\pi+\frac{\pi}{2},(2k+1)\pi+\frac{3\pi}{4}\right)
\end{gathered}
\]
上单调递增，在
\[
\begin{gathered}
\left(2k\pi,2k\pi+\frac{\pi}{4}\right),\quad \left(2k\pi+\frac{\pi}{2},2k\pi+\frac{3\pi}{4}\right),\\
\left(2k\pi+\frac{3\pi}{4},(2k+1)\pi\right),\quad \left((2k+1)\pi+\frac{\pi}{4},(2k+1)\pi+\frac{\pi}{2}\right)
\end{gathered}
\]
上单调递减，其中 \(k\in\Z\).
\end{enumerate}
\end{solution}

\begin{problem}
已知椭圆 \(\Gamma\) 的方程为 \(\frac{x^2}{a^2} + \frac{y^2}{b^2} = 1\)（\(a > b > 0\)），点 \(P\) 的坐标为 \((-a, b)\).

\begin{enumerate}
\item 若直角坐标平面上的点 \(M\)，\(A(0, -b)\)，\(B(a, 0)\) 满足 \(\overrightarrow{PM} = \frac{1}{2} (\overrightarrow{PA} + \overrightarrow{PB})\)，求点 \(M\) 的坐标；
\item 设直线 \(l_1: y = k_1x + p\) 交椭圆 \(\Gamma\) 于 \(C\)，\(D\) 两点，交直线 \(l_2: y = k_2x\) 于点 \(E\). 若 \(k_1 \cdot k_2 = -\frac{b^2}{a^2}\)，证明： \(E\) 为 \(CD\) 的中点；
\item 对于椭圆 \(\Gamma\) 上的点 \(Q(a\cos \theta, b\sin \theta)\) （\(0 < \theta < \pi\)） ，如果椭圆 \(\Gamma\) 上存在不同的两点 \(P_1\)，\(P_2\) 使 \(\overrightarrow{PP_1} + \overrightarrow{PP_2} = \overrightarrow{PQ}\)，写出求作点 \(P_1\)，\(P_2\) 的步骤，并求出使 \(P_1\)，\(P_2\) 存在的 \(\theta\) 满足的条件.
\end{enumerate}
\end{problem}

\begin{solution}
\begin{enumerate}
\item 由 \(P=(-a,b)\)、\(A=(0,-b)\)、\(B=(a,0)\)，得 \(\overrightarrow{PA}=(a,-2b)\)、\(\overrightarrow{PB}=(2a,-b)\). 因而 \(\overrightarrow{PM}=\frac{1}{2}\left(\overrightarrow{PA}+\overrightarrow{PB}\right)=\left(\frac{3a}{2},-\frac{3b}{2}\right)\)，从而 \(M=P+\overrightarrow{PM}=\left(\frac{a}{2},-\frac{b}{2}\right)\).

\item 设 \(C(x_1,y_1)\)、\(D(x_2,y_2)\). 将 \(y=k_1x+p\) 代入椭圆方程，得 \(\left(a^2k_1^2+b^2\right)x^2+2a^2k_1px+a^2p^2-a^2b^2=0\). 由韦达定理，\(x_1+x_2=-\frac{2a^2k_1p}{a^2k_1^2+b^2}\). 因此 \(CD\) 的中点坐标为
\(\left(\frac{x_1+x_2}{2},\frac{y_1+y_2}{2}\right)=\left(-\frac{a^2k_1p}{a^2k_1^2+b^2},\frac{b^2p}{a^2k_1^2+b^2}\right)\)，其中第二个坐标由 \(y_i=k_1x_i+p\) 得到.

由 \(E\in l_1\cap l_2\)，有 \(k_1x_E+p=k_2x_E\)，即 \(x_E=\frac{p}{k_2-k_1}\)、\(y_E=k_2x_E\). 因 \(k_1k_2=-\frac{b^2}{a^2}\) 且 \(k_1\neq0\)，所以 \(k_2=-\frac{b^2}{a^2k_1}\)，从而 \(x_E=-\frac{a^2k_1p}{a^2k_1^2+b^2}\)、\(y_E=\frac{b^2p}{a^2k_1^2+b^2}\). 这与上述中点坐标相同，故 \(E\) 为 \(CD\) 的中点.

\item 设 \(F\) 为 \(PQ\) 的中点. 由 \(P=(-a,b)\)、\(Q=(a\cos\theta,b\sin\theta)\)，得 \(F=\left(\frac{a(\cos\theta-1)}{2},\frac{b(1+\sin\theta)}{2}\right)\). 若 \(P_1\)、\(P_2\) 存在且满足题设向量等式，则 \(F\) 必为 \(P_1P_2\) 的中点，因为 \(\overrightarrow{PP_1}+\overrightarrow{PP_2}=2\overrightarrow{PF}=\overrightarrow{PQ}\).

作直线 \(OF\)，其斜率为 \(k_2=-\frac{b(1+\sin\theta)}{a(1-\cos\theta)}\). 再过 \(F\) 作直线 \(l\)，使其斜率为 \(k_1=-\frac{b^2}{a^2k_2}=\frac{b(1-\cos\theta)}{a(1+\sin\theta)}\). 将直线 \(l\) 与椭圆 \(\Gamma\) 的两个交点作出，记为 \(P_1\)、\(P_2\). 根据第（2）问，\(F\) 是弦 \(P_1P_2\) 的中点；又因为 \(F\) 是 \(PQ\) 的中点，所以这些点满足题设向量等式.

要使交点为不同的两个点，须且只须 \(F\) 在椭圆内部，即 \(\frac{x_F^2}{a^2}+\frac{y_F^2}{b^2}<1\). 代入 \(F\) 的坐标，得 \(\frac{(1-\cos\theta)^2+(1+\sin\theta)^2}{4}<1\)，即 \(\sin\theta-\cos\theta<\frac{1}{2}\). 于是 \(\sqrt{2}\sin\left(\theta-\frac{\pi}{4}\right)<\frac{1}{2}\). 结合 \(0<\theta<\pi\)，可得 \(0<\theta<\frac{\pi}{4}+\arcsin\frac{\sqrt{2}}{4}\).
\end{enumerate}
\end{solution}
